% Glossário em português brasileiro do Volume I.
% IDs e símbolos preservam notations.tex; carregar apenas no driver realanal.tex.

\newglossaryentry{not:def}{
  name={$A \coloneqq B$},
  description={definir $A$ como igual a $B$},
}

\newglossaryentry{not:set}{
  name={$\{ 1,2,3\}$},
  description={conjunto com os elementos indicados},
}

\newglossaryentry{not:equal}{
  name={$A = B$},
  description={$A$ e $B$ são iguais},
}

\newglossaryentry{not:composition}{
  name={$f \circ g$},
  description={composição de funções},
}

\newglossaryentry{not:equivclass}{
  name={$[a]$},
  description={classe de equivalência de $a$},
}

\newglossaryentry{not:cardinality}{
  name={$\abs{A}$},
  description={cardinalidade de um conjunto $A$},
}

\newglossaryentry{not:powerset}{
  name={$\sP(P)$},
  description={conjunto das partes de $A$},
}

\newglossaryentry{not:emptyset}{
  name={$\emptyset$},
  description={conjunto vazio},
}

\newglossaryentry{not:natnums}{
  name={$\N$},
  description={números naturais: $1,2,3,\ldots$},
}

\newglossaryentry{not:ints}{
  name={$\Z$},
  description={números inteiros: $\ldots,-2,-1,0,1,2,\ldots$},
}

\newglossaryentry{not:rationals}{
  name={$\Q$},
  description={números racionais},
}

\newglossaryentry{not:reals}{
  name={$\R$},
  description={números reais},
}

\newglossaryentry{not:complex}{
  name={$\C$},
  description={números complexos},
}

\newglossaryentry{not:euclidspace}{
  name={$\R^n$},
  description={espaço euclidiano de dimensão $n$},
}

\newglossaryentry{not:func}{
  name={$f \colon A \to B$},
  description={função com domínio $A$ e contradomínio $B$},
}

\newglossaryentry{not:dirimage}{
  name={$f(A)$},
  description={imagem direta de $A$ por $f$},
}

\newglossaryentry{not:invimage}{
  name={$f^{-1}(A)$},
  description={imagem inversa de $A$ por $f$},
}

\newglossaryentry{not:invfunction}{
  name={$f^{-1}$},
  description={função inversa},
}

\newglossaryentry{not:abs}{
  name={$\abs{x}$},
  description={valor absoluto},
}

\newglossaryentry{not:closedbdinterval}{
  name={$[a,b]$},
  description={intervalo fechado e limitado},
}

\newglossaryentry{not:openbdinterval}{
  name={$(a,b)$},
  description={intervalo aberto e limitado},
}

\newglossaryentry{not:halfopenbdinterval}{
  name={$(a,b]$, $[a,b)$},
  description={intervalos limitados semiabertos},
}

\newglossaryentry{not:closedubdinterval}{
  name={$[a,\infty)$, $(-\infty,b]$},
  description={intervalos fechados e ilimitados},
}

\newglossaryentry{not:openubdinterval}{
  name={$(a,\infty)$, $(-\infty,b)$},
  description={intervalos abertos e ilimitados},
}

\newglossaryentry{not:sequence}{
  name={$\{x_n\}_{n=1}^\infty$},
  description={sequência},
}

\newglossaryentry{not:subsequence}{
  name={$\{x_{n_i}\}_{i=1}^\infty$},
  description={subsequência},
}

\newglossaryentry{not:sup}{
  name={$\sup\,E$},
  description={supremo de $E$},
}

\newglossaryentry{not:inf}{
  name={$\inf\,E$},
  description={ínfimo de $E$},
}

\newglossaryentry{not:min}{
  name={$\min\,E$},
  description={mínimo de $E$},
}

\newglossaryentry{not:max}{
  name={$\max\,E$},
  description={máximo de $E$},
}

\newglossaryentry{not:supf}{
  name={$\displaystyle \sup_{x\in D} f(x)$},
  description={supremo de $f(D)$},
}

\newglossaryentry{not:inff}{
  name={$\displaystyle \inf_{x\in D} f(x)$},
  description={ínfimo de $f(D)$},
}

\newglossaryentry{not:extreal}{
  name={$\R^*$},
  description={números reais estendidos},
}

\newglossaryentry{not:infinity}{
  name={$\infty$},
  description={infinito},
}

\newglossaryentry{not:limseq}{
  name={$\displaystyle \lim_{n\to\infty} x_n$},
  description={limite de uma sequência},
}

\newglossaryentry{not:limfunc}{
  name={$\displaystyle \lim_{x \to c} f(x)$},
  description={limite de uma função},
}

\newglossaryentry{not:limitasarrows}{
  name={$f(x) \to L$ quando $x \to c$},
  description={$f(x)$ tende a $L$ quando $x$ tende a $c$},
}

\newglossaryentry{not:limfuncinf}{
  name={$\displaystyle \lim_{x \to \infty} f(x)$,
$\displaystyle \lim_{x \to -\infty} f(x)$},
  description={limite de uma função no infinito},
}

\newglossaryentry{not:onesidedlim}{
  name={$\displaystyle \lim_{x \to c^{+}} f(x)$,
$\displaystyle \lim_{x \to c^{-}} f(x)$},
  description={limites laterais de uma função},
}

\newglossaryentry{not:limsupseq}{
  name={$\displaystyle \limsup_{n\to\infty} x_n$},
  description={limite superior},
}

\newglossaryentry{not:liminfseq}{
  name={$\displaystyle \liminf_{n\to\infty} x_n$},
  description={limite inferior},
}

\newglossaryentry{not:finitesum}{
  name={$\displaystyle \sum_{n=1}^k a_n$},
  description={soma $a_1 + a_2 + \cdots+ a_k$},
}

\newglossaryentry{not:series}{
  name={$\displaystyle \sum_{n=1}^\infty a_n$},
  description={série},
}

\newglossaryentry{not:riemannint}{
  name={$\displaystyle \int_a^b f$, $\displaystyle \int_a^b f(x) \, dx$},
  description={integral de Riemann de $f$ em $[a,b]$},
}

\newglossaryentry{not:integrablefunc}{
  name={$\sR\bigl([a,b]\bigr)$},
  description={funções integráveis à Riemann em $[a,b]$},
}

\newglossaryentry{not:setbuilder}{
  name={$\{ x \in S : P(x) \}$},
  description={notação de construção de conjuntos},
}

\newglossaryentry{not:subset}{
  name={$A \subset B$},
  description={$A$ é subconjunto de $B$},
}

\newglossaryentry{not:subsetneq}{
  name={$A \subsetneq B$},
  description={$A$ é subconjunto próprio de $B$},
}

\newglossaryentry{not:eltof}{
  name={$x \in S$},
  description={$x$ é elemento de $S$},
}

\newglossaryentry{not:noteltof}{
  name={$x \notin S$},
  description={$x$ não é elemento de $S$},
}

\newglossaryentry{not:intersection}{
  name={$A \cap B$},
  description={interseção de $A$ e $B$},
}

\newglossaryentry{not:setminus}{
  name={$A \setminus B$},
  description={diferença entre $A$ e $B$, formada pelos elementos de $A$ que não pertencem a $B$},
}

\newglossaryentry{not:complement}{
  name={$B^c$},
  description={complemento de $B$, formado pelos elementos que não pertencem a $B$},
}

\newglossaryentry{not:cartprod}{
  name={$A \times B$},
  description={produto cartesiano de $A$ e $B$},
}

\newglossaryentry{not:union}{
  name={$A \cup B$},
  description={união de $A$ e $B$},
}

\newglossaryentry{not:infunion}{
  name={$\displaystyle \bigcup_{n=1}^\infty A_n$},
  description={união dos conjuntos $A_n$, para todos os $n \in \N$},
}

\newglossaryentry{not:arbitunion}{
  name={$\displaystyle \bigcup_{\lambda \in I} A_\lambda$},
  description={união dos conjuntos $A_\lambda$, para todos os $\lambda \in I$},
}

\newglossaryentry{not:infintersection}{
  name={$\displaystyle \bigcap_{n=1}^\infty A_n$},
  description={interseção dos conjuntos $A_n$, para todos os $n \in \N$},
}

\newglossaryentry{not:arbitintersection}{
  % Intervenção editorial transparente: a fonte exibe \bigcup, mas esta entrada e o texto definem uma interseção; adotamos \bigcap.
  name={$\displaystyle \bigcap_{\lambda \in I} A_\lambda$},
  description={interseção dos conjuntos $A_\lambda$, para todos os $\lambda \in I$},
}

\newglossaryentry{not:lowerdarbouxsum}{
  name={$L(P,f)$},
  description={soma inferior de Darboux de $f$ sobre a partição $P$},
}

\newglossaryentry{not:upperdarbouxsum}{
  name={$U(P,f)$},
  description={soma superior de Darboux de $f$ sobre a partição $P$},
}

\newglossaryentry{not:lowerdarboux}{
  name={$\displaystyle \underline{\int_a^b}f$},
  description={integral inferior de Darboux},
}

\newglossaryentry{not:upperdarboux}{
  name={$\displaystyle \overline{\int_a^b}f$},
  description={integral superior de Darboux},
}

\newglossaryentry{not:exp}{
  name={$\exp(x)$, $e^x$},
  description={função exponencial},
}

\newglossaryentry{not:ln}{
  name={$\ln(x)$, $\log(x)$},
  description={função logaritmo natural},
}

\newglossaryentry{not:pow}{
  name={$x^y$},
  description={potência de base $x>0$ e expoente $y\in\R$},
}

\newglossaryentry{not:e}{
  name={$e$},
  description={número de Euler, base do logaritmo natural},
}

\newglossaryentry{not:contfuncs}{
  name={$C(S,\R)$},
  description={funções contínuas $f \colon S \to \R$},
}

\newglossaryentry{not:contdifffuncs}{
  name={$C^1(S,\R)$},
  description={funções continuamente diferenciáveis $f \colon S \to \R$},
}

\newglossaryentry{not:derivative}{
  name={$f'(x)$, $\frac{df}{dx}$, $\frac{d}{dx}\bigl( f(x) \bigr)$},
  description={derivada de $f$},
}

\newglossaryentry{not:secondthirdfourthder}{
  name={$f''$, $f'''$, $f''''$},
  description={derivadas de segunda, terceira e quarta ordem de $f$},
}

\newglossaryentry{not:nthder}{
  name={$f^{(n)}$},
  description={derivada de ordem $n$ de $f$},
}

\newglossaryentry{not:ordereq}{
  name={$x = y$},
  description={$x$ é igual a $y$},
}

\newglossaryentry{not:orderlt}{
  name={$x < y$},
  description={$x$ é menor que $y$},
}

\newglossaryentry{not:orderleq}{
  name={$x \leq y$},
  description={$x$ é menor ou igual a $y$},
}

\newglossaryentry{not:ordergt}{
  name={$x > y$},
  description={$x$ é maior que $y$},
}

\newglossaryentry{not:ordergeq}{
  name={$x \geq y$},
  description={$x$ é maior ou igual a $y$},
}

\newglossaryentry{not:uniformnorm}{
  name={$\norm{f}_S$},
  description={norma uniforme de $f$ em $S$},
}

\newglossaryentry{not:diam}{
  name={$\operatorname{diam}(S)$},
  description={diâmetro de $S$},
}

\newglossaryentry{not:openball}{
  name={$B(p,\delta)$, $B_X(p,\delta)$},
  description={bola aberta em um espaço métrico},
}

\newglossaryentry{not:closedball}{
  name={$C(p,\delta)$, $C_X(p,\delta)$},
  description={bola fechada em um espaço métrico},
}

\newglossaryentry{not:closure}{
  %FIXME: this is weird, we need protect
  name={$\protect\widebar{A}$},
  description={fecho de $A$},
}

\newglossaryentry{not:interior}{
  name={$A^\circ$},
  description={interior de $A$},
}

\newglossaryentry{not:boundary}{
  name={$\partial A$},
  description={fronteira de $A$},
}
